% diversity_reduction_conditions.tex [NEW]
% Conditions for diversity reduction, anti-pole failure mode, and
% practical argument for why the condition is generically satisfied.
%
% Intended to follow the local (diversity_reduction_theory.tex) and
% global (diversity_reduction_geometric.tex) sections.
%
% === NOTES ===
%
% This section arose from noticing that neither the Taylor nor the
% geometric result actually proves diversity REDUCTION relative to the
% unsteered baseline in full generality. The Taylor result diverges
% when m = mu + s approx 0, and the geometric result bounds absolute spread
% but doesn't compare to the unsteered case. This file states the
% necessary condition, shows what goes wrong at the anti-pole, and
% argues that the condition is generically satisfied in practice.
%
% Citation keys used:
%   \citet{turntrout2023norms} -- residual stream norm growth
%
% === END NOTES ===

\subsection{When Does Steering Reduce Diversity?}
\label{sec:conditions}

The results of Sections~\ref{sec:covariance} and~\ref{sec:geometric-bound}
establish that the post-normalization distribution concentrates around the
pole $\hat{s}$ with spread shrinking as $O(1/\|s\|)$.
However, neither result directly compares the steered diversity to the
\emph{unsteered} baseline.
In this section, we identify the condition under which steering genuinely
reduces diversity relative to the unsteered model, show that the condition
can fail, and argue that it is satisfied in all practical scenarios.

\subsubsection{The Condition}

Consider the unsteered case ($s = 0$).
The normalization map sends $x$ to $\hat{x} = x / \|x\|$, projecting
the distribution $\mathcal{D}$ onto the sphere.
The total variance of the unsteered output is governed (via the
first-order analysis) by $1/\|\mu\|^2$, and the steered output by
$1/\|\mu + s\|^2$.
Steered diversity is lower than unsteered diversity when the denominator
grows, i.e., when:
%
\begin{equation}
  \|\mu + s\| > \|\mu\|.
  \label{eq:diversity-condition}
\end{equation}
%
Squaring both sides and expanding:
%
\begin{equation}
  \|\mu\|^2 + 2\mu^\top s + \|s\|^2 > \|\mu\|^2,
\end{equation}
%
which simplifies to:
%
\begin{equation}
  2\mu^\top s + \|s\|^2 > 0
  \qquad\Longleftrightarrow\qquad
  \cos\angle(\mu, s) > -\frac{\|s\|}{2\|\mu\|}.
  \label{eq:angle-condition}
\end{equation}
%
This is a weak requirement.
For any $\|s\| > 0$, the right-hand side is negative, so the condition
is satisfied whenever $\cos\angle(\mu, s) \geq 0$
(i.e., the steering vector is in the same half-space as the mean),
and even permits obtuse angles up to
$\arccos(-\|s\| / 2\|\mu\|)$.
Only when the steering vector is nearly anti-parallel to the mean and
small relative to it can the condition fail.

\subsubsection{Failure at the Anti-Pole}
\label{sec:anti-pole}

When $s \approx -\mu$, the shifted mean $m = \mu + s$ is near the origin.
In this regime:
%
\begin{itemize}
  \item The first-order approximation breaks down:
    the condition $\|\delta\| \ll \|m\|$ is violated, and the predicted
    variance $\mathrm{tr}(\Sigma_x) / \|m\|^2$ diverges.
  \item Geometrically, the distribution of shifted vectors $\{x + s\}$
    straddles the origin.
    Points on opposite sides of the origin are projected to opposite sides
    of the sphere, so the normalization map acts as a
    \emph{magnifying glass}: small Euclidean differences become large
    angular differences.
  \item The global concentration bound
    (Proposition~\ref{prop:concentration}) becomes vacuous: when
    $\|s\| \approx \|\mu\|$, the bound
    $2\,\mathbb{E}[\|x\|] / \|s\| \approx 2$ saturates the maximum
    possible chordal distance.
\end{itemize}
%
This is not merely a limitation of our proof technique.
The diversity genuinely increases in this regime: the normalization map
near the origin has unbounded Lipschitz constant
(the Jacobian $J_g(v) = (I - \hat{v}\hat{v}^\top)/\|v\|$ has operator
norm $1/\|v\|$), so it amplifies input variation.

\subsubsection{Why the Condition Is Generically Satisfied}
\label{sec:generic}

We argue that the anti-pole regime is not encountered in practice, for
two independent reasons.

\paragraph{Geometric argument in high dimensions.}
In $\mathbb{R}^d$ with large $d$, two structured vectors are unlikely to
be nearly anti-parallel.
Two independent random unit vectors have cosine similarity with mean $0$
and standard deviation $1/\sqrt{d}$; for $d = 4096$, achieving
$\cos\angle(\mu, s) < -1/64$ (which is still far from the threshold
$-\|s\|/2\|\mu\|$) would be a multi-sigma event.
While $\mu$ and $s$ are not random---they are structured vectors
reflecting the model's learned representations---the same concentration
of measure makes near-anti-alignment a priori unlikely without
deliberate adversarial construction.

\paragraph{Mechanistic argument.}
Steering vectors are typically computed as activation \emph{differences}
between prompts exhibiting contrasting behaviors
(e.g., eval-aware minus eval-unaware).
These differences encode a behavioral contrast direction.
The mean residual stream $\mu$ at a given layer encodes average content:
positional information, syntactic structure, and semantic features
aggregated over the input distribution.
There is no mechanistic reason for a behavioral contrast direction to
point toward the negation of the average content.
To construct $s \approx -\mu$, one would need a steering vector that
cancels the aggregate content of the residual stream---a pathological
intervention that would destroy all model computation, not just reduce
diversity.

\paragraph{Norm growth.}
\citet{turntrout2023norms} document that residual stream norms grow
exponentially with layer depth (approximately $1.045\times$ per layer
in GPT-2-XL).
Steering vectors, being fixed additive perturbations, do not participate
in this learned growth schedule.
At deeper layers, $\|\mu\|$ is therefore typically much larger than
$\|s\|$, making the condition \eqref{eq:angle-condition} even easier
to satisfy: the threshold $-\|s\|/2\|\mu\|$ approaches zero, and
essentially any non-adversarial steering direction suffices.

\subsubsection{Recommended Treatment}

For the formal results, we state the condition
$\|\mu + s\| > \|\mu\|$ as an explicit assumption.
For all experimental results, we verify it empirically by computing
$\|\mu + s\| / \|\mu\|$ at the injection layer and confirming the
ratio exceeds 1.
This is a lightweight check (requiring only the mean residual stream
norm, which is already available from the evaluation pipeline) and
provides a per-experiment certificate that the diversity-reduction
theory applies.
